% GENERATED FILE. Edit canonical lessons, then run npm run generate:books.
% canonical-source-hash: fnv1a64:e67afb99b794c1af
% canonical-lessons: MW-C60-khad, MW-C60-lun, MW-C60-mirach, MW-C60-dahi, MW-C60-tel

\chapter{Sugar, Salt, Chilli, Yoghurt, Oil}
\label{ch:mw-sugar-salt-chilli-yoghurt-oil}
\hlchaptermodality{60}

\begin{chapteropening}
\textbf{By the end of this chapter:} \emph{I can say sugar, salt, a chilli, yoghurt and oil.}

Complete the last lesson of chapter 60: I can say sugar, salt, a chilli, yoghurt and oil.
\end{chapteropening}

\section[\texorpdfstring{\mw{खांड} (khā̃ḍ)}{खांड (khā̃ḍ)}]{\mw{खांड} (khā̃ḍ) — sugar}

\label{lesson:MW-C60-khad}



\begin{quote}
Before the new one: say the Marwadi for a tomato, then the Marwadi for a carrot.
\end{quote}

\subsection*{You'll want to know: \mw{खांड}}
\textbf{\mw{खांड}} — \emph{khā̃ḍ} — \textquotedblleft{}sugar\textquotedblright{}.

\begin{cousinweb}
\textbf{\mw{खांड}} (\emph{khā̃ḍ}) is a tadbhava, a word inherited from Sanskrit and reshaped by sound change. It comes from \emph{khaṇḍa}, \textquotedblleft{}partially dried molasses\textquotedblright{}. CDIAL lists Old Marwari \emph{khāḍa} under it (CDIAL 3800).
\end{cousinweb}

\begin{grammarlens}[title={What you've built}]
One more word for this chapter.
\end{grammarlens}

\subsection*{Your turn}
\begin{itemize}
\raggedright
  \item \emph{Say it:} \emph{khā̃ḍ}
  \item \emph{Say it:} \emph{khā̃ḍ}, once more
  \item \emph{Say it:} say \emph{khā̃ḍ}
  \item \emph{Recall:} say the Marwadi for a tomato, then the Marwadi for a carrot, then say \emph{khā̃ḍ} again
\end{itemize}

\begin{tcolorbox}[breakable,colback=teal!4,colframe=teal!35!black,title={Before you move on}]
What does \mw{खांड} mean? (\textquotedblleft{}Sugar\textquotedblright{}.) Say it once more.
\end{tcolorbox}
\section[\texorpdfstring{\mw{लूण} (lūṇ)}{लूण (lūṇ)}]{\mw{लूण} (lūṇ) — salt}

\label{lesson:MW-C60-lun}



\begin{quote}
Before the new one: say the Marwadi for a carrot, then the Marwadi for sugar.
\end{quote}

\subsection*{You'll want to know: \mw{लूण}}
\textbf{\mw{लूण}} — \emph{lūṇ} — \textquotedblleft{}salt\textquotedblright{}.

\begin{cousinweb}
\textbf{\mw{लूण}} (\emph{lūṇ}) is a tadbhava, a word inherited from Sanskrit and reshaped by sound change. It comes from \emph{lavaṇa}, \textquotedblleft{}salt\textquotedblright{}. CDIAL lists Marwari \emph{lūṇ} under it (CDIAL 10978). CDIAL prefers a Middle Indo-Aryan blend of \emph{lavana} with \emph{lūna}, \textquotedblleft{}a piece of rock-salt\textquotedblright{}, for the long \emph{ū}.
\end{cousinweb}

\begin{grammarlens}[title={What you've built}]
One more word for this chapter.
\end{grammarlens}

\subsection*{Your turn}
\begin{itemize}
\raggedright
  \item \emph{Say it:} \emph{lūṇ}
  \item \emph{Say it:} \emph{lūṇ}, once more
  \item \emph{Say it:} say \emph{lūṇ}
  \item \emph{Recall:} say the Marwadi for a carrot, then the Marwadi for sugar, then say \emph{lūṇ} again
\end{itemize}

\begin{tcolorbox}[breakable,colback=teal!4,colframe=teal!35!black,title={Before you move on}]
What does \mw{लूण} mean? (\textquotedblleft{}Salt\textquotedblright{}.) Say it once more.
\end{tcolorbox}
\section[\texorpdfstring{\mw{मिरच} (mirach)}{मिरच (mirach)}]{\mw{मिरच} (mirach) — a chilli}

\label{lesson:MW-C60-mirach}



\begin{quote}
Before the new one: say the Marwadi for sugar, then the Marwadi for salt.
\end{quote}

\subsection*{You'll want to know: \mw{मिरच}}
\textbf{\mw{मिरच}} — \emph{mirach} — \textquotedblleft{}a chilli\textquotedblright{}.

\begin{cousinweb}
\textbf{\mw{मिरच}} (\emph{mirach}) is a tadbhava, a word inherited from Sanskrit and reshaped by sound change. It comes from \emph{marīca}, \textquotedblleft{}peppercorn\textquotedblright{}. CDIAL gives no Marwari form in this entry (CDIAL 9875); the match is through Hindi \emph{mirc}. CDIAL takes the Sanskrit word itself from an Austroasiatic language. The name moved from black pepper to the chilli.
\end{cousinweb}

\begin{grammarlens}[title={What you've built}]
One more word for this chapter.
\end{grammarlens}

\subsection*{Your turn}
\begin{itemize}
\raggedright
  \item \emph{Say it:} \emph{mirach}
  \item \emph{Say it:} \emph{mirach}, once more
  \item \emph{Say it:} say \emph{mirach}
  \item \emph{Recall:} say the Marwadi for sugar, then the Marwadi for salt, then say \emph{mirach} again
\end{itemize}

\begin{tcolorbox}[breakable,colback=teal!4,colframe=teal!35!black,title={Before you move on}]
What does \mw{मिरच} mean? (\textquotedblleft{}A chilli\textquotedblright{}.) Say it once more.
\end{tcolorbox}
\section[\texorpdfstring{\mw{दही} (dahī)}{दही (dahī)}]{\mw{दही} (dahī) — yoghurt}

\label{lesson:MW-C60-dahi}



\begin{quote}
Before the new one: say the Marwadi for salt, then the Marwadi for a chilli.
\end{quote}

\subsection*{You'll want to know: \mw{दही}}
\textbf{\mw{दही}} — \emph{dahī} — \textquotedblleft{}yoghurt\textquotedblright{}.

\begin{cousinweb}
\textbf{\mw{दही}} (\emph{dahī}) is a tadbhava, a word inherited from Sanskrit and reshaped by sound change. It comes from \emph{dadhi}, \textquotedblleft{}thick sour milk\textquotedblright{}. CDIAL gives no Marwari form in this entry (CDIAL 6146); the match is through Hindi \emph{dahī}.
\end{cousinweb}

\begin{grammarlens}[title={What you've built}]
One more word for this chapter.
\end{grammarlens}

\subsection*{Your turn}
\begin{itemize}
\raggedright
  \item \emph{Say it:} \emph{dahī}
  \item \emph{Say it:} \emph{dahī}, once more
  \item \emph{Say it:} say \emph{dahī}
  \item \emph{Recall:} say the Marwadi for salt, then the Marwadi for a chilli, then say \emph{dahī} again
\end{itemize}

\begin{tcolorbox}[breakable,colback=teal!4,colframe=teal!35!black,title={Before you move on}]
What does \mw{दही} mean? (\textquotedblleft{}Yoghurt\textquotedblright{}.) Say it once more.
\end{tcolorbox}
\section[\texorpdfstring{\mw{तेल} (tel)}{तेल (tel)}]{\mw{तेल} (tel) — oil}

\label{lesson:MW-C60-tel}



\begin{quote}
Before the new one: say the Marwadi for a chilli, then the Marwadi for yoghurt.
\end{quote}

\subsection*{You'll want to know: \mw{तेल}}
\textbf{\mw{तेल}} — \emph{tel} — \textquotedblleft{}oil\textquotedblright{}.

\begin{cousinweb}
\textbf{\mw{तेल}} (\emph{tel}) is a tadbhava, a word inherited from Sanskrit and reshaped by sound change. It comes from \emph{taila}, \textquotedblleft{}sesamum oil, oil\textquotedblright{}. CDIAL gives no Marwari form in this entry (CDIAL 5958); the match is through Hindi \emph{tel} and Gujarati \emph{tel}. It is built on \emph{tila}, \textquotedblleft{}sesame\textquotedblright{}.
\end{cousinweb}

\begin{grammarlens}[title={What you've built}]
One more word for this chapter.
\end{grammarlens}

\subsection*{Your turn}
\begin{itemize}
\raggedright
  \item \emph{Say it:} \emph{tel}
  \item \emph{Say it:} \emph{tel}, once more
  \item \emph{Say it:} say \emph{tel}
  \item \emph{Recall:} say the Marwadi for a chilli, then the Marwadi for yoghurt, then say \emph{tel} again
\end{itemize}

\begin{tcolorbox}[breakable,colback=teal!4,colframe=teal!35!black,title={Before you move on}]
What does \mw{तेल} mean? (\textquotedblleft{}Oil\textquotedblright{}.) Say it once more.
\end{tcolorbox}
